%===== File containing the back cover of the document =====%
%                                                          %
% Copyright (C) ISI - All Rights Reserved                  %
% Proprietary                                              %
% Written by Med Hossam <med.hossam@gmail.com>, April 2016 %
%                                                          %
% @author: HEDHILI Med Houssemeddine                       %
% @linkedin: http://tn.linkedin.com/in/medhossam           %
%==========================================================%

%== It's advised to not modify the content of this file ===%
% To set your information, go to global_config.tex file    %
%==========================================================%

\thispagestyle{backcover}
\newgeometry{bottom=25mm,left=15mm,top=20mm,right=15mm}

\begin{changemargin}{3mm}{0cm}
    \begin{minipage}[c]{0.96\columnwidth}
        
        
        % {\ifthenelse{\boolean{wantToTypeCompanyAddress}}
        % {% IF TRUE
        %     \vskip5mm
        % }{\vskip8mm}}
        
        \selectlanguage{french}
        
        {\LARGE\textbf{Résumé}}
        \vskip1mm
            \begingroup
                \large
                Le projet « Supervision automatisée par l’IA et collecte intelligente d’informations », réalisé au sein de la société SerOps dans le cadre de l’obtention du diplôme national d’ingénieur en informatique, vise à centraliser et analyser les journaux et métriques des systèmes en temps réel grâce à un modèle d’intelligence artificielle. L’application permet la détection automatique des anomalies, la génération de recommandations et l’amélioration de la gestion des incidents. L’interface utilisateur a été développée avec React, le backend avec Django, et les données sont collectées et traitées via Elasticsearch et Fluentd, puis stockées dans PostgreSQL. Cette solution offre une supervision efficace et sécurisée des systèmes, facilitant la prise de décision et le suivi des performances..
            \endgroup
        \vskip1mm
        {\textbf{Mots clés : }
            \begingroup
                IA, Supervision, Logs, Métriques, React, Django, Elasticsearch, Fluentd, PostgreSQL
            \endgroup
        }
        
        {\vskip25mm}
        
        \selectlanguage{english}
        {\LARGE\textbf{Abstract}}
        \vskip1mm
            \begingroup
                \large
                The project “Automated Supervision with AI and Intelligent Information Collection”, carried out at SerOps as part of the completion of the National Engineering Diploma in Computer Science, aims to centralize and analyze system logs and metrics in real time using an artificial intelligence model. The application enables automatic anomaly detection, generates recommendations, and improves incident management. The user interface was developed with React, the backend with Django, and data is collected and processed using Elasticsearch and Fluentd, then stored in PostgreSQL. This solution provides effective and secure system supervision, facilitating decision-making and performance monitoring.
            \endgroup
        \vskip1mm
        {\textbf{Keywords : }
            \begingroup
              AI, Supervision, Logs, Metrics, React, Django, Elasticsearch, Fluentd, PostgreSQL.
            \endgroup
        }
    \end{minipage}
    
\end{changemargin}